\begin{center}\begin{tikzpicture}[font=\small,>=stealth,line join=round]
\draw[gray!40,rounded corners=3pt] (0,0) rectangle (14.2,8.3);\node at (7.1,7.700000000000001) {例6：条件保证周期4，曲线却每隔2已经重复};
\begin{scope}[shift={(1.6882352941176473,3.7409090909090907)},x=1.1764705882352942cm,y=1.3939393939393938cm]\draw[->,gray] (-0.5,0)--(4.6,0) node[right]{$x$};\draw[->,gray] (0,-1.5)--(0,1.8) node[above]{$y$};\draw[blue,thick,solid,domain=-0.2:4.3,samples=180] plot(\x,{cos(deg(pi*\x))});\fill[orange!85!black] (0,1) circle(2.2pt);\fill[orange!85!black] (2,1) circle(2.2pt);\fill[orange!85!black] (4,1) circle(2.2pt);\end{scope}\node[anchor=west,align=left,text width=5.0cm] at (8.4,4.0) {$f=\cos(\pi x)$\\[9pt]$f(x+2)=f(x-2)$\\[9pt]推出4是一个周期\\[9pt]最小正周期为2\\[9pt]图中4含两个完整重复};
\node[text width=12.6cm,align=center] at (7.1,.65) {只有进一步使用这条具体函数的性质，才能确定最小正周期。};
\end{tikzpicture}\end{center}
